\documentclass[12pt]{article}
\usepackage{fullpage}
\usepackage{amsmath,amssymb,amsthm}

\newtheorem{lemma}{Lemma}
\newtheorem{proposition}{Proposition}

\DeclareMathOperator{\Lk}{Lk}
\newcommand{\bbarotimes}{\mathbin{\overline{\otimes}}}
\newcommand{\cA}{\mathcal A}
\newcommand{\cH}{\mathcal H}
\newcommand{\C}{\mathbb C}

\begin{document}

\title{Proper proximality of irreducible graph-product von Neumann algebras}
\author{}
\date{}
\maketitle

\section*{Problem statement and interpretation}

I use the standard graph-product meaning of irreducible: the finite graph
\(\Gamma=(V,E)\) is not a non-trivial join, equivalently the complement
\(\Gamma^c\) is connected.  The traces are faithful normal traces.  For
\(U\subset V\), \(M_U\) denotes the graph product over the induced subgraph and
\(E_U:M\to M_U\) its canonical expectation.

Put \(A_N=\mathbb B(L^2N)\) and
\[
        \cA_N=(A_N^{\sharp_J})^*
\]
for the simultaneous \(N\)- and \(JNJ\)-normal bidual of
Ding--Kunnawalkam Elayavalli--Peterson (DKEP) \cite{DKEP}.  If \(X\) is an
\(N\)-boundary piece, \(q_X\in\cA_N\) denotes the support projection of the
corner \((K_X(N)^{\sharp_J})^*\).  Thus
\[
 \widetilde S_X(N)=
 \{T\in\cA_N:[T,a]\in q_X\cA_Nq_X\hbox{ for all }a\in JNJ\}.
\]
For \(X=\mathbb K(L^2N)\) write \(q_0\) and \(\widetilde S(N)\).  We use DKEP
Lemma 8.5 for the trivial group action: if \(N\) is a separable finite factor,
then proper proximality relative to \(X\) is equivalent to the absence of an
\(N\)-central state on \(\widetilde S_X(N)\) restricting to the trace.

\section*{Normal supports and graph-product boundary corners}

We first record the support calculus needed below.  Let \(A=\mathbb B(L^2N)\),
viewed as an \(N\)- and \(JNJ\)-\(C^*\)-algebra.  If \(B\subset A\) is a
hereditary \(C^*\)-subalgebra invariant under the multiplier actions of \(N\)
and \(JNJ\), let \(s_B\in A^{**}\) be its ordinary open support.  In each one-normal category, the paragraph before DKEP Remark 2.8
identifies the normal bidual of \(B\) with the cutdown
\(s_B(A^{\sharp})^*s_B\) inside \(A^{**}\).  DKEP Proposition 2.10 identifies
the simultaneous normal biduals with the intersections of the two one-normal
biduals; hence
\[
        (B^{\sharp_J})^*
        = \cA_N\cap s_BA^{**}s_B                                      \tag{1}
\]
inside \(A^{**}\).  This is a corner of \(\cA_N\); denote its unit by \(q(B)\).
Thus \(T\in q(B)\cA_Nq(B)\) exactly when \(T\in\cA_N\) and
\(T=s_BT s_B\) in \(A^{**}\).  No centrality of the normal-support projection is
being used here.

For a DKEP boundary piece \(X\), this \(q(X)\) is the same projection as
\(q_X\).  Indeed, \(q_X=q(K_X(N))\) by definition.  If
\(\phi\in A^{\sharp_J}\) vanishes on \(X\), then DKEP Theorem 5.6 with
\(X=Y\) gives \(\phi(T)=0\) for every
\(T\in K_X^{\infty,1}(N)\); in particular \(\phi\) vanishes on \(K_X(N)\).
The converse is trivial because \(X\subset K_X(N)\).  Hence \(X\) and
\(K_X(N)\) have the same annihilator in \(A^{\sharp_J}\), and therefore the
same weak-* closed corner in \(\cA_N\).

The following consequence will be used repeatedly.  If \(B,C,D\) are such
hereditary subalgebras and
\[
        s_C(s_BA^{**}s_B)s_C\subset s_DA^{**}s_D,                  \tag{2}
\]
then
\[
        q(C)\bigl(q(B)\cA_Nq(B)\bigr)q(C)\subset q(D)\cA_Nq(D).     \tag{3}
\]
This follows at once from (1): the left-hand side is contained in
\(\cA_N\cap s_DA^{**}s_D\).

For \(U\subset V\), set
\[
        \cH_U=L^2M\otimes_{M_U}L^2M .
\]
If \(\Omega_U=1\otimes1\), DKEP's coefficient operator is
\(T_{\Omega_U}=e_{M_U}\), since \(T_{\Omega_U}\widehat x=\widehat{E_U(x)}\).
Thus DKEP Example 5.11 identifies the coefficient boundary of \(\cH_U\) with
the subalgebra boundary \(X_{M_U}\) of Example 3.4.  If \(Z\subset U\), then
\(e_{M_Z}\le e_{M_U}\), whence \(X_{M_Z}\subset X_{M_U}\).

\begin{lemma}[boundary-corner intersection]\label{lem:intersection}
Let \(q_U=q_{X_{M_U}}\).  For all \(U,W\subset V\),
\[
        q_W(q_U\cA_Mq_U)q_W\subset q_{U\cap W}\cA_Mq_{U\cap W}.   \tag{4}
\]
Consequently, if \(U_1,\ldots,U_m\subset V\) and \(I=\cap_iU_i\), then
\[
R=q_{U_m}q_{U_{m-1}}\cdots q_{U_1}q_{U_2}\cdots q_{U_{m-1}}q_{U_m}
\in q_I\cA_Mq_I .
\]
\end{lemma}

\begin{proof}
Let \(s_U\) be the ordinary support of \(X_{M_U}\) in
\(A_M^{**}\), and put \(I=U\cap W\).  By the preceding support calculus, it
suffices to prove
\[
        s_W(s_UA_M^{**}s_U)s_W\subset s_IA_M^{**}s_I.             \tag{5}
\]

The boundary \(X_{M_U}\) has a contractive approximate unit whose elements are
norm limits of finite sums of coefficients of finite Connes fusions of
\(\cH_U\) and \(\overline{\cH_U}\).  Indeed, \(X_{M_U}\) is hereditarily
generated by positive coefficients \(T_\xi^*T_\xi\); the usual net
\(f_\varepsilon(\sum_{j=1}^nT_{\xi_j}^*T_{\xi_j})\), with
\(f_\varepsilon(0)=0\), is such an approximate unit, and polynomial
approximation plus the identities \(T_{\bar\xi}=T_\xi^*\) and
\(T_{\xi\otimes\eta}=T_\eta T_\xi\) rewrites the monomials as finite-fusion
coefficients.

Choose such approximate units \(h_\alpha^W\in X_{M_W}\) and
\(h_\beta^U\in X_{M_U}\).  We claim
\[
        h_\alpha^Wh_\beta^U\in s_IA_M^{**}s_U,\qquad
        h_\beta^Uh_\alpha^W\in s_UA_M^{**}s_I .                  \tag{6}
\]
It is enough to check this on finite-fusion coefficients.  If
\(a=T_{\xi_1}^*T_{\xi_2}\) is a coefficient of a fusion all of whose labels lie
in \(W\), and \(b=T_{\eta_1}^*T_{\eta_2}\) is one with labels in \(U\), then
\[
        ab
        =T_{\xi_1\otimes\overline{\xi_2}\otimes\eta_1}^{\,*}T_{\eta_2}.
                                                                        \tag{7}
\]
CDHJKN identify \(\cH_U\) with their correspondence \(H_U(V,V)\), and their
fusion rule \cite[Theorem 5.4 and Proposition 5.5]{CDHJKN}, applied
repeatedly, decomposes the correspondence containing
\(\xi_1\otimes\overline{\xi_2}\otimes\eta_1\) as a direct sum of
\(\cH_Z\)'s with \(Z\subset U\cap W\).  Since \(X_{M_Z}\subset X_{M_I}\) for
such \(Z\), and off-diagonal coefficients in a direct sum are controlled by the
standard \(2\times2\) positivity argument, its diagonal coefficient is supported
by \(s_I\).  Also \(T_{\eta_2}^*T_{\eta_2}\) is supported by \(s_U\).
The elementary rectangular support fact
\(T^*T\in sA^{**}s,\ S^*S\in tA^{**}t\Rightarrow T^*S\in sA^{**}t\)
gives the first inclusion in (6); the second is the adjoint.

Now let \(x\in s_UA_M^{**}s_U\).  Since \(h_\alpha^W\to s_W\) and
\(h_\beta^U\to s_U\) weak-* in \(A_M^{**}\),
\[
h_\alpha^Wh_\beta^U\,x\,h_\gamma^Uh_\delta^W\in s_IA_M^{**}s_I
\]
by (6), and successive weak-* limits yield (5).  Applying (3) gives (4).  The
palindromic assertion follows by induction from (4).
\end{proof}

\section*{A finite boundary criterion}

\begin{lemma}\label{lem:finite}
Let \(N\) be a separable non-amenable finite factor.  Let \(X_1,\ldots,X_m\) be
boundary pieces and put \(q_i=q_{X_i}\).  Assume
\[
        q_mq_{m-1}\cdots q_1q_2\cdots q_{m-1}q_m\in q_0\cA_Nq_0.  \tag{8}
\]
If \(N\) is properly proximal relative to each \(X_i\), then \(N\) is properly
proximal.
\end{lemma}

\begin{proof}
Assume \(N\) is not properly proximal.  DKEP Lemma 8.5 gives an \(N\)-central
state \(\omega\) on \(\widetilde S(N)\) with \(\omega|_N=\tau\).  Fix \(i\) and
write \(q=q_i\).  Since \(q\) commutes with \(N\) and \(JNJ\),
\(q^\perp\widetilde S_{X_i}(N)q^\perp\subset\widetilde S(N)\).  If
\(\omega(q^\perp)>0\), compression by \(q^\perp\) gives an \(N\)-central state
on \(\widetilde S_{X_i}(N)\) whose restriction to \(N_+\) is dominated by a
scalar multiple of \(\tau\), hence is normal; because \(N\) is a factor this
restriction is \(\tau\).  This contradicts relative proper proximality, so
\(\omega(q_i^\perp)=0\) for every \(i\).

Extend \(\omega\) to a state on \(\cA_N\).  In the GNS representation every
\(q_i\) fixes the cyclic vector, so the positive contraction in (8) has
\(\omega\)-value \(1\).  Since it lies in the corner \(q_0\cA_Nq_0\), we get
\(\omega(q_0)=1\).  DKEP's canonical \(N\)- and \(JNJ\)-bimodular u.c.p. map
\(\Theta:\mathbb B(L^2N)\to q_0\cA_Nq_0\) sends \(1\) to \(q_0\) and is the
identity on \(N\) after compression.  Thus \(\omega\circ\Theta\) is a Connes
hypertrace on \(\mathbb B(L^2N)\), contradicting non-amenability of \(N\).
\end{proof}

\section*{A relative Bass--Serre paradox}

\begin{proposition}\label{prop:BS}
Let \(N=P*_{B}Q\) with trace-preserving expectations and suppose
\(Q=B\bbarotimes A\).  Assume that \(A\) contains a trace-zero unitary \(a\) and
that \(P\) contains unitaries \(w_1,w_2\) satisfying
\[
        E_B(w_i)=0,\qquad E_B(w_1^*w_2)=0 .
\]
Then \(N\) is properly proximal relative to \(X_B\).
\end{proposition}

\begin{proof}
Let \(H_P=L^2P\ominus L^2B\) and \(H_Q=L^2Q\ominus L^2B\).  In the AFP
decomposition of \(L^2N\), let \(p\) and \(q\) be the projections onto reduced
words whose first letter lies in \(H_P\) and \(H_Q\), respectively; then
\(1=e_B+p+q\).  Direct inspection of reduced words gives, for
\(R_x=Jx^*J\),
\[
[p,R_x]=R_xe_B-e_BR_x,\ [q,R_x]=0\quad(x\in P\ominus B),
\]
and the same formula with \(p,q\) interchanged for \(x\in Q\ominus B\); elements
of \(B\) commute with \(p,q\).  These commutators are \(N\)- and \(JNJ\)-translates
of \(e_B\), so DKEP Lemma 6.1 implies \(p,q\in S_{X_B}(N)\).

Let \(\varphi\) be an \(N\)-central state on \(S_{X_B}(N)\) normal on \(N\).
Left multiplication by a centered \(P\)-letter sends \(qL^2N\) into \(pL^2N\),
and \(q w_1^*w_2q=0\) because \(E_B(w_1^*w_2)=0\).  Hence
\(w_1qw_1^*+w_2qw_2^*\le p\).  Similarly \(apa^*\le q\).  Centrality gives
\(2\varphi(q)\le\varphi(p)\le\varphi(q)\), so \(\varphi(p)=\varphi(q)=0\).

The unitary \(h=w_1a\) is \(B\)-Haar: all non-zero powers are reduced alternating
words, hence have \(B\)-expectation zero.  The projections
\(h^ne_Bh^{-n}\) are pairwise orthogonal and have the same \(\varphi\)-value,
so \(\varphi(e_B)=0\).  This contradicts \(1=e_B+p+q\).  DKEP Theorem 6.2
therefore gives proper proximality relative to \(X_B\).
\end{proof}

\section*{Application to graph products}

Fix \(v\in V\) and let \(L_v=\Lk_\Gamma(v)\).  The graph product splits as
\[
        M\cong M_{V\setminus\{v\}}*_{M_{L_v}}
        (M_{L_v}\bbarotimes M_v).                                \tag{9}
\]
Let \(C(v)=\{r\ne v:r\not\sim_\Gamma v\}\).  Since \(\Gamma^c\) is connected,
\(C(v)\ne\emptyset\).  If \(C(v)\) contains distinct \(r,t\), choose trace-zero
unitaries \(u_r\in M_r,u_t\in M_t\) and put \(w_1=u_r,w_2=u_t\).  If
\(C(v)=\{r\}\), connectedness and \(|V|\ge3\) give
\(s\in L_v\) with \(s\not\sim_\Gamma r\); put \(w_1=u_r,w_2=u_su_r\).  In both
cases reduced-word calculus gives
\[
        E_{M_{L_v}}(w_i)=0,\qquad E_{M_{L_v}}(w_1^*w_2)=0 .
\]
The trace-zero unitary in \(M_v\) is the unitary \(a\) in (9), so
Proposition \ref{prop:BS} shows that \(M\) is properly proximal relative to
\(X_{M_{L_v}}\) for every \(v\in V\).

Enumerate \(V=\{v_1,\ldots,v_m\}\).  Since no vertex belongs to its own link,
\(\cap_iL_{v_i}=\emptyset\).  Lemma \ref{lem:intersection} gives
\[
 q_{L_{v_m}}\cdots q_{L_{v_1}}q_{L_{v_2}}\cdots q_{L_{v_m}}
 \in q_{\emptyset}\cA_Mq_{\emptyset}.
\]
Here \(M_\emptyset=\C\), and \(X_{M_\emptyset}=\mathbb K(L^2M)\).

Finally, CDHJKN Main Theorem 0.5 implies that \(M\) is a factor: the possible
obstructions would give either a universal vertex of \(\Gamma\) or an isolated
exceptional non-edge component, both impossible because \(\Gamma^c\) is
connected with at least three vertices.  CDHJKN Proposition 6.3 similarly
excludes amenability, since amenability would force the non-edge components of
\(\Gamma^c\) to be isolated exceptional components.  Thus \(M\) is a separable
non-amenable finite factor.  Lemma \ref{lem:finite} applies and proves that
\(M\) is properly proximal.

\begin{thebibliography}{9}

\bibitem{DKEP}
C.~Ding, S.~Kunnawalkam Elayavalli, and J.~Peterson,
\emph{Properly proximal von Neumann algebras},
Duke Math. J. \textbf{172} (2023), 2821--2894.

\bibitem{CDHJKN}
I.~Charlesworth, R.~de Santiago, B.~Hayes, D.~Jekel,
S.~Kunnawalkam Elayavalli, and B.~Nelson,
\emph{On the structure of graph product von Neumann algebras},
Publ. Res. Inst. Math. Sci. \textbf{61} (2025), 713--762.

\end{thebibliography}

\end{document}
